\subsection{Native Run on a Large Trunk Model}
\label{sec:spn-10b-bruno}

This section records what was needed to run the static optimization of
Section~\ref{sec:spn-10b} on a trunk model with 446 muscles, 105 locked
coordinates and no coupling constraints (issue \#12761). Only dimensionless
ratios and wall times are given; the licensed model and its numbers stay out
of the repository, and its hash is unpinned, so such a receipt is
\texttt{unqualified} by design.

\paragraph{Muscle Model Policy.}
OpenSim~4.6 does not implement \texttt{ignore\_tendon\_compliance} for
\texttt{Thelen2003Muscle}, so the rigid-tendon evaluation fails with an opaque
error. The request option \texttt{muscle\_model\_policy} is \texttt{native}
(default; the stage fails with a message naming the option) or
\texttt{millard\_equivalent}, which rewrites the muscles as
\texttt{Millard2012EquilibriumMuscle} carrying the same maximal force, fibre
and tendon lengths and pennation. This is a modelling substitution: the
muscle force--length--velocity curves differ. The policy enters the request
hash when non-default, and the receipt records the swapped count and note.

\paragraph{Locked Coordinates.}
A locked coordinate ignores \texttt{setValue}, is held at zero acceleration in
forward dynamics, has no moment arm, and is ignored by the inverse-dynamics
solver. A drive that moves such a coordinate therefore makes the
inverse-dynamics and joint-reaction accelerations disagree (a relative
mismatch of order $10^{-2}$ against $10^{-6}$ unlocked). The default is to
fail the drive stage and list the offending coordinates (truncated).
\texttt{unlock\_driven\_coordinates} unlocks exactly those coordinates; the
count and names are recorded. Holding them constant instead changes the load
distribution and is left as a follow-up (ADR-0055).

\paragraph{Solver and Screen.}
SLSQP did not converge at 446 muscles. For 100 or more muscles stage two uses
Clarabel: one quadratic program for $p=2$ and iteratively reweighted quadratic
programs for $p=3$ (weights $w_m=a_m^{p-2}$, floor $10^{-6}$, tolerance
$10^{-10}$). A path-length central-difference screen (step $10^{-4}$ rad,
tolerance $10^{-13}$ m) drops muscle--coordinate pairs without a moment arm;
locked coordinates are always evaluated natively. A test requires screened and
unscreened evaluation to agree. The screen threshold and the solver
(\texttt{slsqp}, \texttt{clarabel\_qp}, \texttt{clarabel\_irls}) are recorded
in the stage provenance.

\paragraph{Observed Behaviour and Limits.}
Both $p=2$ and $p=3$ studies on a synthetic swing finished in about one
minute with status \texttt{unqualified}: the asset hash is unpinned, the
gravity-only pelvis residual exceeds its budget, and reserve shares were
flagged on every event. The Fujii coupling law covers the lumbar levels only,
so on a T1--L5 model the drive fails; the regional law
\texttt{bruno2015\_regional} with \texttt{restrict\_law\_to\_model} is used
instead, and a lumbar-only sensitivity option is not provided. In an informal
comparison of that law with the uniform reference, compression in the lumbar
and lower thoracic levels varied least; moments, shear and upper-thoracic
loads did not. No scientific acceptance follows from these runs.

\paragraph{Reproduction.}
Set \texttt{UD\_SPINE\_BRUNO\_SOURCE} to the model file, build a
\texttt{SpineStudyRequest} with \texttt{asset\_id}, the regional law,
\texttt{muscle\_model\_policy="millard\_equivalent"} and
\texttt{unlock\_driven\_coordinates=True}, and call \texttt{run\_spine\_study}.
